\subsubsection{Resultados de la encuesta a alumnos}
\label{sec:resultados-alumnos}

La Tabla~\ref{tab:resultados-alumnos} resume las respuestas a las siete preguntas con escala de caritas. Seis de los 17 alumnos respondieron \textit{Sí} en todas ellas.

\begin{table}[H]
\centering
\caption{Respuestas de los alumnos a las preguntas con escala de caritas ($n=17$).}
\label{tab:resultados-alumnos}
\renewcommand{\arraystretch}{1.3}
\begin{tabular}{|p{6.4cm}|c|c|c|c|}
\hline
\textbf{Pregunta} & \textbf{Sí} & \textbf{Más o menos} & \textbf{No} & \textbf{Dos marcas} \\ \hline
1. ¿Te gustó Graf-IA? & 16 (94\,\%) & 0 & 0 & 1 \\ \hline
2. ¿Entendiste lo que tenías que hacer? & 12 (71\,\%) & 3 & 0 & 2 \\ \hline
3. ¿Fue fácil escribir en la pantalla? & 10 (59\,\%) & 3 & 1 & 3 \\ \hline
4. ¿Pudiste leer bien las letras de la pantalla? & 13 (76\,\%) & 1 & 1 & 2 \\ \hline
5. ¿Te ayudó escuchar las instrucciones? & 15 (88\,\%) & 1 & 0 & 1 \\ \hline
6. ¿Los dibujos te ayudaron a saber qué escribir? & 16 (94\,\%) & 0 & 0 & 1 \\ \hline
7. ¿Te gustaría practicar tus letras con Graf-IA otra vez? & 16 (94\,\%) & 1 & 0 & 0 \\ \hline
\end{tabular}
\end{table}

Las preguntas sobre el agrado general, la utilidad de los dibujos y la disposición a volver a practicar obtuvieron la mayor proporción de respuestas afirmativas (16 de 17). La facilidad para escribir en la pantalla recibió 10 respuestas \textit{Sí}, la cifra más baja de las siete preguntas. En esa pregunta hubo tres respuestas \textit{Más o menos}, una respuesta \textit{No} y tres encuestas con dos marcas. La pregunta sobre la comprensión de la actividad también recibió tres respuestas \textit{Más o menos}.

La Tabla~\ref{tab:preferencias-alumnos} muestra las respuestas a las preguntas de opción única.

\begin{table}[H]
\centering
\caption{Preferencias de los alumnos ($n=17$).}
\label{tab:preferencias-alumnos}
\renewcommand{\arraystretch}{1.3}
\begin{tabular}{|p{5.2cm}|p{5.2cm}|c|}
\hline
\textbf{Pregunta} & \textbf{Opción} & \textbf{Alumnos} \\ \hline
\multirow{3}{=}{8. ¿Dónde te gusta más escribir?} & En los dos (tableta y cuaderno) & 11 \\ \cline{2-3}
 & En la tableta & 4 \\ \cline{2-3}
 & En mi cuaderno & 2 \\ \hline
\multirow{4}{=}{9. ¿Qué fue lo que más te gustó?} & Escribir con el dedo & 8 \\ \cline{2-3}
 & Ver mi avance & 7 \\ \cline{2-3}
 & Escuchar & 1 \\ \cline{2-3}
 & Los dibujos & 1 \\ \hline
\end{tabular}
\end{table}

Once de los 17 alumnos prefirieron combinar la escritura en la tableta con la del cuaderno. Escribir con el dedo y ver el avance fueron las funciones que más les gustaron.

En la pregunta abierta sobre lo más difícil, 14 alumnos escribieron la palabra \textit{ejercicio}, con variantes como \textit{ejevcicio} o \textit{ejercio}, y tres respondieron con un dibujo sin texto. Las respuestas escritas se refieren al ejercicio en general y dejan abierta la pregunta sobre qué parte causó dificultad.
